\section{Практическая часть}
\noindent\textbf{Замечание.}
Заметим, что в данной работе рассматривается 4 разных длины фактически одной волны:
\begin{enumerate}
	\item $\lambda^{\text{т}}_{\text{в}}$ - длина волны в волноводе, рассчитанная теоретически
	\item $\lambda^{\text{т}}_{\text{с}}$ - длина волны в свободном пространстве, рассчитанная теоретически
	\item $\lambda^{\text{п}}_{\text{в}}$ - длина волны в волноводе, измеренная на практике
	\item $\lambda^{\text{п}}_{\text{с}}$ - длина волны в свободном пространстве, измеренная на практике
\end{enumerate}


\noindent\textbf{Начальные условия.}
В работе используется прямоугольный волновод размерами $23 \times 10$~мм. Шаг перемещения зонда и антенны при снятии экспериментальных зависимостей составляет $3$~мм. Расстояние от антенны до отражающего щита
\begin{equation}
	X = 283~\text{см} = 2{,}83~\text{м}.
\end{equation}

\noindent\textbf{Проверка применимости метода.}
Щит находится в зоне Фраунгофера:
\begin{equation}
	X \gg \frac{\max\{l_1^2, l_2^2\}}{{\lambda^{\text{п}}_{\text{с}}}}
	= \frac{179{,}56~\text{см}^2}{3{,}32~\text{см}}
	\approx 54~\text{см},
\end{equation}
то есть условие выполняется.

\subsection{Задание 1. Поле при отсутствии принимаемого сигнала}

Перед антенной помещается щит с поглощающим покрытием ($\Gamma \approx 0$). Снята зависимость интенсивности поля $|E|^2$ от координаты зонда $x$ (см.~\cref{fig:graph:1}, \cref{tab:practice1}).

При отсутствии отражённой компоненты уравнение \cref{eq:th:13} упрощается:
\begin{equation}
	|E|^2 \approx 1 + 2\Gamma_\kappa \cos(2hx + \varphi_\kappa).
	\label{eq:pr:1.1}
\end{equation}

Максимум интенсивности достигается при $\cos(2hx + \varphi_\kappa) = 1$, а минимум при $\cos(2hx + \varphi_\kappa) = -1$. Для нормированного поля:
\begin{align*}
	|E|^2_{max}& = 1 + 2\Gamma_\kappa \\
	|E|^2_{min}& = 1 - 2\Gamma_\kappa.
\end{align*}


Отсюда можно получить
\begin{equation}
	\Gamma_\kappa
	= \frac{|E|^2_{\max} - |E|^2_{\min}}{2(|E|^2_{\max} + |E|^2_{\min})}
	= \frac{54 - 45}{2(54 + 45)}
	\approx 0{,}045.
	\label{eq:pr:1.2}
\end{equation}



\noindent\textbf{Длина волны.}
По \cref{tab:practice1} максимумы показаний расположены при $x = 0$, $24$ и $48$~мм. Фаза в \cref{eq:pr:1.1} изменяется как $2hx$, где $h = 2\pi/{\lambda^{\text{п}}_{\text{в}}}$, поэтому расстояние между соседними максимумами равно половине длины волны в волноводе:
\begin{equation}
	{\lambda^{\text{п}}_{\text{в}}} =2 \times 24~\text{мм} = 4{,}8~\text{см}.
\end{equation}

Связь длины волны в свободном пространстве с ${\lambda^{\text{п}}_{\text{в}}}$:
\begin{equation}
	{\lambda^{\text{п}}_{\text{в}}} = \frac{{\lambda^{\text{п}}_{\text{с}}}}{\sqrt{1 - ({\lambda^{\text{п}}_{\text{с}}}/2a)^2}},
	\qquad\Longrightarrow\qquad
	{\lambda^{\text{п}}_{\text{с}}} = \left(\frac{1}{\bigl(\lambda^{\text{п}}_{\text{в}}\bigr)^2} + \frac{1}{(2a)^2}\right)^{-1/2},
\end{equation}
где $a = 2{,}3$~см~--- размер широкой стенки волновода; рассматривается основная мода $H_{10}$. По измеренной длине волны получаем
\begin{equation}
	{\lambda^{\text{п}}_{\text{с}}} = \left(\frac{1}{4{,}8^2} + \frac{1}{4{,}6^2}\right)^{-1/2}
	\approx 3{,}32~\text{см} = 0{,}0332~\text{м}.
\end{equation}
При указанной частоте генератора $f = 9{,}0$~ГГц длина волны в свободном пространстве
\begin{equation}
	{\lambda^{\text{т}}_{\text{с}}} = \frac{c}{f}
	= \frac{3 \cdot 10^8}{9{,}0 \cdot 10^9}
	\approx 3{,}33~\text{см}.
\end{equation}
Расхождение с экспериментальной оценкой ($3{,}32$~см) составляет около $0{,}4\%$ и объясняется дискретностью шага ($3$~мм), неточностью определения координат экстремумов и погрешностью отсчёта. В дальнейших экспериментальных расчётах используется ${\lambda^{\text{п}}_{\text{с}}} = 0{,}0332$~м.

\noindent\textbf{Коэффициент отражения от конца тракта и КПД.}
Так как детектор квадратичный, коэффициент бегущей волны вычисляется через корень из отношения показаний:
\begin{equation}\label{eq:pr:3.1}
	\kappa_{\text{к}} = \sqrt{\frac{I_{\min}}{I_{\max}}}
	= \sqrt{\frac{45}{54}}
	\approx 0{,}9129,
\end{equation}
\begin{equation}\label{eq:pr:3.2}
	\Gamma_{\text{к}} = \frac{1 - \kappa_{\text{к}}}{1 + \kappa_{\text{к}}}
	= \frac{1 - 0{,}9129}{1 + 0{,}9129}
	\approx 0{,}0455.
\end{equation}
Оценка КПД системы по рассогласованию. Мощность пропорциональна квадрату амплитуды поля: $P \propto |E|^2$. В данном случае $P_{\text{отпр}} = P_{\text{подан}} - P_{\text{отр}}$.
\begin{align}
	\eta_{\text{согл}} &= \frac{P_{\text{отпр}}}{P_{\text{подан}}}
	= \frac{P_{\text{подан}} - P_{\text{отр}}}{P_{\text{подан}}}
	= 1 - \frac{P_{\text{отр}}}{P_{\text{подан}}} \notag\\
	&= 1 - \frac{\Gamma_{\text{к}}^2 P_{\text{подан}}}{P_{\text{подан}}}
	= 1 - \Gamma_{\text{к}}^2.
\end{align}

\begin{equation}
	\eta_{\text{согл}} = 1 - \Gamma_{\text{к}}^2
	= 1 - 0{,}0455^2
	\approx 0{,}9979 = 99{,}79\%.
\end{equation}

Для второго задания фиксируется координата, при которой $\cos(2hx + \varphi_{\text{к}}) = 0$. Между соседними экстремумами $x = 24$~мм и $x = 36$~мм:
\begin{equation}
	X_0 = \frac{24 + 36}{2} = 30~\text{мм}.
\end{equation}

\subsection{Задание 2. Определение КНД по изменению расстояния до щита}

После удаления поглощающего щита антенна перемещалась относительно металлического щита, измерительная линия оставалась в $X_0 = 30$~мм. Результаты~--- в \cref{tab:practice2}, зависимость~--- на \cref{fig:graph:2}.

Максимальное и минимальное показания детектора:
\begin{equation}
	I_{\max} = 51~\text{мА}, \qquad I_{\min} = 47~\text{мА}.
\end{equation}
Коэффициент отражения от металлического щита получим подставив \cref{eq:pr:3.1} в \cref{eq:pr:3.2}:
\begin{equation}
	\Gamma_2 = \frac{\sqrt{I_{\max}} - \sqrt{I_{\min}}}{\sqrt{I_{\max}} + \sqrt{I_{\min}}}
	= \frac{\sqrt{51} - \sqrt{47}}{\sqrt{51} + \sqrt{47}}
	\approx 0{,}0204.
	\label{eq:pr:2.1}
\end{equation}
По формуле \cref{eq:th:14}, используя измеренную длину волны ${\lambda^{\text{п}}_{\text{с}}} = 0{,}0332$~м:
\begin{equation}
	D_2 = \frac{8\pi X}{{\lambda^{\text{п}}_{\text{с}}}}\Gamma_2
	= \frac{8\pi \cdot 2{,}83}{0{,}0332} \cdot 0{,}0204
	\approx 43{,}7,
	\qquad
	D_2 \approx 16{,}4~\text{дБ}.
\end{equation}

Аппроксимация измерений функцией $I(\Delta X) = A + B\cos(4\pi\Delta X/{\lambda^{\text{п}}_{\text{с}}} + \varphi)$ даёт ${\lambda^{\text{п}}_{\text{с},2}} \approx 34{,}0$~мм, что согласуется с оценкой по частоте генератора (${\lambda^{\text{т}}_{\text{с}}} \approx 33{,}3$~мм) в пределах около $2\%$.

\noindent\textbf{Проверка малости.}
Для применимости линейного приближения:
\begin{equation}
	\Gamma_{\text{к}}^2 \approx 0{,}0021, \qquad
	\Gamma_2^2 \approx 0{,}0004, \qquad
	\Gamma_{\text{к}}\Gamma_2 \approx 0{,}0009.
\end{equation}
Все три величины малы по сравнению с единицей.

\subsection{Задание 3. Определение КНД по максимуму $\widetilde{\Gamma}$}

Для каждой позиции антенны по $E_{\max}$ и $E_{\min}$ (\cref{tab:practice3}) рассчитывались
\begin{equation}
	\kappa(\Delta X) = \sqrt{\frac{E_{\min}}{E_{\max}}},
	\qquad
	\widetilde{\Gamma}(\Delta X) = \frac{1 - \kappa}{1 + \kappa}.
\end{equation}
Зависимость~--- на \cref{fig:graph:3}.

Максимум достигается при $\Delta X = 60$~мм:
\begin{equation}
	\widetilde{\Gamma}_{\max} \approx 0{,}0627.
\end{equation}
С учётом $\widetilde{\Gamma}_{\max} \approx \Gamma + \Gamma_{\text{к}}$:
\begin{equation}
	\Gamma_3 = 0{,}0627 - 0{,}0455 \approx 0{,}0172,
\end{equation}
\begin{equation}
	D_3 = \frac{8\pi X}{{\lambda^{\text{п}}_{\text{с}}}}\Gamma_3
	= \frac{8\pi \cdot 2{,}83}{0{,}0332} \cdot 0{,}0172
	\approx 36{,}8,
	\qquad
	D_3 \approx 15{,}7~\text{дБ}.
\end{equation}

\noindent\textbf{Проверка малости.}
\begin{equation}
	\Gamma_{\text{к}}^2 \approx 0{,}0021, \qquad
	\Gamma_3^2 \approx 0{,}0003, \qquad
	\Gamma_{\text{к}}\Gamma_3 \approx 0{,}0008.
\end{equation}
Все три величины малы, приближение корректно.

\subsection{Задание 4. Теоретический расчёт КНД}

Расчёт по интегралу Кирхгофа--Гюйгенса (см.~\cref{eq:th:1.1,eq:th:1.2,eq:th:1.3,eq:th:1.4,eq:th:1.5,eq:th:1.6}) при предположении о синфазности и равномерности поля на раскрыве, с округлённой длиной волны ${\lambda^{\text{п}}_{\text{с}}} \approx 3{,}3$~см, даёт
\begin{equation}
	D_{\text{теор}} = \frac{4\pi S}{\bigl(\lambda^{\text{п}}_{\text{с}}\bigr)^2}
	= \frac{4\pi \cdot 9{,}2 \cdot 13{,}4}{3{,}3^2}
	\approx 142{,}3.
\end{equation}
В действительности поле на раскрыве рупора несинфазно: фаза меняется от центра к краям из-за сферичности фронта волны, что уменьшает эффективную площадь раскрыва и КНД по сравнению с \cref{eq:th:1.6}. Реальный КНД $D_{\text{реал}} = g \cdot 4\pi S/\bigl(\lambda^{\text{п}}_{\text{с}}\bigr)^2$ с апертурным коэффициентом $g < 1$.

\subsection{Задание 5. Сравнение эксперимента и теории}

Сводка результатов (КНД в децибелах рассчитывается как $D_{\text{дБ}} = 10\log_{10}D$):

\begin{center}
	\begin{tabular}{lcc}
		\toprule
		Метод & $D$ & $D$, дБ \\
		\midrule
		Задание 2 (по $\Gamma_2$) & $43{,}7$ & $16{,}4$ \\
		Задание 3 (по $\widetilde{\Gamma}_{\max}$) & $36{,}8$ & $15{,}7$ \\
		Теория \cref{eq:th:1.6} & $142{,}3$ & $21{,}5$ \\
		\bottomrule
	\end{tabular}
\end{center}

Расхождение двух экспериментальных методов:
\begin{equation}
	\delta_{\text{мет}} = \frac{|D_2 - D_3|}{(D_2 + D_3)/2} \cdot 100\%
	= \frac{|43{,}7 - 36{,}8|}{(43{,}7 + 36{,}8)/2} \cdot 100\%
	\approx 17{,}1\%.
\end{equation}

Расхождение эксперимента с теорией:
\begin{equation}
	\delta_{\text{теор}} = \frac{|D_{\text{теор}} - (D_2 + D_3)/2|}{D_{\text{теор}}} \cdot 100\%
	= \frac{|142{,}3 - 40{,}25|}{142{,}3} \cdot 100\%
	\approx 71{,}7\%.
\end{equation}

\noindent\textbf{Причины расхождений.}
\begin{itemize}
	\item Основная причина расхождения с теорией~--- неоднородность фазы поля на раскрыве рупора: формула $D = 4\pi S/\bigl(\lambda^{\text{п}}_{\text{с}}\bigr)^2$ предполагает синфазное и равномерное поле. В реальном рупоре фаза квадратично нарастает от центра к краям, что снижает КНД. По полученным результатам апертурный коэффициент составляет $g \approx 40{,}25/142{,}3 \approx 0{,}28$.
	\par
	\item Расхождение между двумя экспериментальными методами связано с дискретностью шага ($3$~мм), неточностью определения экстремумов $E_{\min}$, $E_{\max}$, остаточным отражением от элементов установки и приближением, в котором пренебрегается квадратичными по $\Gamma$ членами.
	\item Дополнительный вклад дают неидеальное согласование тракта и конечные размеры отражающего щита.
\end{itemize}
